% GENERATED FILE. Edit canonical lessons, then run npm run generate:books.
% canonical-source-hash: fnv1a64:97601b3a649232a4
% canonical-lessons: PA-C222-fut-return, PA-C222-karanga, PA-C222-karoge, PA-C222-karanga-review, PA-C222-ki-nahin-future, PA-C222-payoff

\chapter{What Will You Do Tomorrow?}
\label{ch:pa-a2-karanga}
\hlchaptermodality{222}

\begin{chapteropening}
\textbf{By the end of this chapter:} \emph{I can say and ask what people will do tomorrow, and say what I will not do.}

Say what you, others and a group will do with karanga, karega, karoge and karange, and ask or deny it.
\end{chapteropening}

\section[\texorpdfstring{\pa{ਜਾਵਾਂਗਾ} · \pa{ਜਾਵੇਗਾ} … (jāvāṅgā · jāvegā …)}{ਜਾਵਾਂਗਾ · ਜਾਵੇਗਾ … (jāvāṅgā · jāvegā …)}]{Return to going tomorrow}

\label{lesson:PA-C222-fut-return}



\begin{quote}
\emph{Hear:} \emph{jāvāṅgā; jāvegā; jāvāṅge} Say \textbf{\pa{ਜਾਵਾਂਗਾ} · \pa{ਜਾਵੇਗਾ} · \pa{ਜਾਵਾਂਗੇ}}.
\end{quote}

\begin{grammarlens}[title={Grammar Lens}]
Those three futures needed \textbf{\pa{ਵ}} after the vowel of \textbf{\pa{ਜਾ}}. The verb \textbf{\pa{ਕਰਨਾ}}, \textquotedblleft{}to do\textquotedblright{}, has a stem that ends in a consonant, \textbf{\pa{ਕਰ}}. Keep it ready.
\end{grammarlens}

\subsection*{Your turn}
\begin{itemize}
\raggedright
  \item \emph{Hear:} \emph{I will go tomorrow} Say it with the ending that fits you.
  \item \emph{Hear:} \emph{she will go tomorrow} Say \textbf{\pa{ਉਹ} \pa{ਕੱਲ੍ਹ} \pa{ਜਾਵੇਗੀ}}.
  \item \emph{Hear:} \emph{to do} Say \textbf{\pa{ਕਰਨਾ}}.
\end{itemize}

\begin{tcolorbox}[breakable,colback=teal!4,colframe=teal!35!black,title={Before you move on}]
Say \emph{we will go}, then \emph{to do}.
\end{tcolorbox}
\section[\texorpdfstring{\pa{ਕਰਾਂਗਾ} · \pa{ਕਰੇਗਾ} (karāṅgā · karegā)}{ਕਰਾਂਗਾ · ਕਰੇਗਾ (karāṅgā · karegā)}]{\pa{ਕਰਾਂਗਾ}, \pa{ਕਰੇਗਾ} — will do}

\label{lesson:PA-C222-karanga}



\begin{quote}
Say \emph{I will go} and \emph{he will go}. Then \emph{to do}.
\end{quote}

\begin{grammarlens}[title={Grammar Lens: the ending straight on}]
A consonant stem takes the endings with no \textbf{\pa{ਵ}} between:

\begin{quote}
\textbf{\pa{ਕਰ}} + \textbf{\pa{ਾਂਗਾ}} $\to$ \textbf{\pa{ਕਰਾਂਗਾ}} · \textbf{\pa{ਕਰ}} + \textbf{\pa{ੇਗਾ}} $\to$ \textbf{\pa{ਕਰੇਗਾ}}
\end{quote}

\begin{quote}
\textbf{\pa{ਮੈਂ} \pa{ਕੱਲ੍ਹ} \pa{ਕੰਮ} \pa{ਕਰਾਂਗਾ।}} — \emph{maiṁ kallh kamm karāṅgā.} — I will work tomorrow (a man speaking).
\end{quote}

\begin{quote}
\textbf{\pa{ਮੈਂ} \pa{ਕੱਲ੍ਹ} \pa{ਕੰਮ} \pa{ਕਰਾਂਗੀ।}} — \emph{maiṁ kallh kamm karāṅgī.} — I will work tomorrow (a woman speaking).
\end{quote}

\begin{quote}
\textbf{\pa{ਅਮਨ} \pa{ਕੰਮ} \pa{ਕਰੇਗਾ।}} — \emph{Aman kamm karegā.} — Aman will work.
\end{quote}

\textbf{\pa{ਕਰੇਗੀ}} is for a woman, and \textbf{\pa{ਕਰਾਂਗੇ}} for \textbf{\pa{ਅਸੀਂ}}.
\end{grammarlens}

\subsection*{Your turn}
\begin{itemize}
\raggedright
  \item \emph{Hear:} \emph{she will work tomorrow} Say \textbf{\pa{ਉਹ} \pa{ਕੱਲ੍ਹ} \pa{ਕੰਮ} \pa{ਕਰੇਗੀ।}}
  \item \emph{Hear:} \emph{we will work} Say \textbf{\pa{ਅਸੀਂ} \pa{ਕੰਮ} \pa{ਕਰਾਂਗੇ।}}
  \item \emph{Read it:} \textbf{\pa{ਕਰਾਂਗਾ}} Point to the stem \textbf{\pa{ਕਰ}}, then the ending.
  \item \emph{Write it:} copy \textbf{\pa{ਕਰਾਂਗਾ}} and \textbf{\pa{ਕਰੇਗਾ}} once
\end{itemize}

\begin{tcolorbox}[breakable,colback=teal!4,colframe=teal!35!black,title={Before you move on}]
\emph{I will do}? (\textbf{\pa{ਕਰਾਂਗਾ}} or \textbf{\pa{ਕਰਾਂਗੀ}}.) \emph{He will do}? (\textbf{\pa{ਕਰੇਗਾ}}.)
\end{tcolorbox}
\section[\texorpdfstring{\pa{ਕਰੋਗੇ} · \pa{ਕਰਨਗੇ} (karoge · karange)}{ਕਰੋਗੇ · ਕਰਨਗੇ (karoge · karange)}]{\pa{ਕਰੋਗੇ}, \pa{ਕਰਨਗੇ} — you will do, they will do}

\label{lesson:PA-C222-karoge}



\begin{quote}
Say \emph{I will work tomorrow} and \emph{we will work}. Then \textbf{\pa{ਤੁਸੀਂ}}.
\end{quote}

\begin{grammarlens}[title={Grammar Lens: two more cells}]
With \textbf{\pa{ਤੁਸੀਂ}} the future of \textbf{\pa{ਕਰਨਾ}} is \textbf{\pa{ਕਰੋਗੇ}}; with \textbf{\pa{ਉਹ}}, for a group, it is \textbf{\pa{ਕਰਨਗੇ}}:

\begin{quote}
\textbf{\pa{ਤੁਸੀਂ} \pa{ਕੱਲ੍ਹ} \pa{ਕੀ} \pa{ਕਰੋਗੇ}?} — \emph{tusīṁ kallh kī karoge?} — What will you do tomorrow?
\end{quote}

\begin{quote}
\textbf{\pa{ਉਹ} \pa{ਕੱਲ੍ਹ} \pa{ਕੰਮ} \pa{ਕਰਨਗੇ।}} — \emph{uh kallh kamm karange.} — They will work tomorrow.
\end{quote}

Here \textbf{\pa{ਕੀ}} is \textquotedblleft{}what\textquotedblright{}: it stands where the answer goes.
\end{grammarlens}

\subsection*{Your turn}
\begin{itemize}
\raggedright
  \item \emph{Hear:} \emph{what will you do tomorrow?} Say \textbf{\pa{ਤੁਸੀਂ} \pa{ਕੱਲ੍ਹ} \pa{ਕੀ} \pa{ਕਰੋਗੇ}?}
  \item \emph{Hear:} \emph{I will go to the market} Answer it with the ending that fits you.
  \item \emph{Read it:} \textbf{\pa{ਉਹ} \pa{ਕੱਲ੍ਹ} \pa{ਕੰਮ} \pa{ਕਰਨਗੇ।}} Give its meaning.
  \item \emph{Write it:} copy \textbf{\pa{ਕਰੋਗੇ}} once
\end{itemize}

\begin{tcolorbox}[breakable,colback=teal!4,colframe=teal!35!black,title={Before you move on}]
\emph{You will do}? (\textbf{\pa{ਕਰੋਗੇ}}.) \emph{They will do}? (\textbf{\pa{ਕਰਨਗੇ}}.)
\end{tcolorbox}
\section[\texorpdfstring{\pa{ਕਰਾਂਗਾ} · \pa{ਕਰੇਗਾ} … (karāṅgā · karegā …)}{ਕਰਾਂਗਾ · ਕਰੇਗਾ … (karāṅgā · karegā …)}]{Hear and say will do again}

\label{lesson:PA-C222-karanga-review}



\begin{quote}
\emph{Hear:} \emph{karāṅgā; karegā; karoge; karange} Say all four.
\end{quote}

\subsection*{Your turn}
\begin{itemize}
\raggedright
  \item \emph{Hear:} \emph{Aman will work tomorrow} Say \textbf{\pa{ਅਮਨ} \pa{ਕੱਲ੍ਹ} \pa{ਕੰਮ} \pa{ਕਰੇਗਾ।}}
  \item \emph{Hear:} \emph{what will you do?} Say \textbf{\pa{ਤੁਸੀਂ} \pa{ਕੀ} \pa{ਕਰੋਗੇ}?}
  \item \emph{Hear:} \emph{we will go home} Say \textbf{\pa{ਅਸੀਂ} \pa{ਘਰ} \pa{ਜਾਵਾਂਗੇ।}}
  \item \emph{Read it:} \textbf{\pa{ਉਹ} \pa{ਕੰਮ} \pa{ਕਰਨਗੇ।}} Give its meaning.
  \item \emph{Write it:} cover the model and write \textbf{\pa{ਮੈਂ} \pa{ਕੰਮ} \pa{ਕਰਾਂਗਾ।}} or its woman's form
\end{itemize}

\begin{tcolorbox}[breakable,colback=teal!4,colframe=teal!35!black,title={Before you move on}]
Say what you will do tomorrow.
\end{tcolorbox}
\section[\texorpdfstring{\pa{ਕੀ} \pa{ਤੁਸੀਂ} \pa{ਕੱਲ੍ਹ} \pa{ਕੰਮ} \pa{ਕਰੋਗੇ}? (kī tusīṁ kallh …)}{ਕੀ ਤੁਸੀਂ ਕੱਲ੍ਹ ਕੰਮ ਕਰੋਗੇ? (kī tusīṁ kallh …)}]{\pa{ਕੀ} \pa{ਤੁਸੀਂ} \pa{ਕੱਲ੍ਹ} \pa{ਕੰਮ} \pa{ਕਰੋਗੇ}? — will you work tomorrow?}

\label{lesson:PA-C222-ki-nahin-future}



\begin{quote}
Say \emph{what will you do tomorrow?} Then \emph{are you a traveller?}
\end{quote}

\begin{grammarlens}[title={Grammar Lens: \pa{ਕੀ} and \pa{ਨਹੀਂ} with the future}]
Nothing moves. \textbf{\pa{ਕੀ}} goes in front to ask yes or no; \textbf{\pa{ਨਹੀਂ}} goes before the verb to say no:

\begin{quote}
\textbf{\pa{ਕੀ} \pa{ਤੁਸੀਂ} \pa{ਕੱਲ੍ਹ} \pa{ਕੰਮ} \pa{ਕਰੋਗੇ}?} — \emph{kī tusīṁ kallh kamm karoge?} — Will you work tomorrow?
\end{quote}

\begin{quote}
\textbf{\pa{ਨਹੀਂ} \pa{ਜੀ}, \pa{ਮੈਂ} \pa{ਕੱਲ੍ਹ} \pa{ਕੰਮ} \pa{ਨਹੀਂ} \pa{ਕਰਾਂਗੀ।}} — \emph{nahīṁ jī, maiṁ kallh kamm nahīṁ karāṅgī.} — No, I won't work tomorrow.
\end{quote}

\begin{quote}
\textbf{\pa{ਕੀ} \pa{ਅਮਨ} \pa{ਕੱਲ੍ਹ} \pa{ਆਵੇਗਾ}?} — \emph{kī Aman kallh āvegā?} — Will Aman come tomorrow?
\end{quote}
\end{grammarlens}

\subsection*{Your turn}
\begin{itemize}
\raggedright
  \item \emph{Hear:} \emph{will she go to school tomorrow?} Say \textbf{\pa{ਕੀ} \pa{ਉਹ} \pa{ਕੱਲ੍ਹ} \pa{ਸਕੂਲ} \pa{ਜਾਵੇਗੀ}?}
  \item \emph{Hear:} \emph{no, she won't go} Say \textbf{\pa{ਨਹੀਂ} \pa{ਜੀ}, \pa{ਉਹ} \pa{ਨਹੀਂ} \pa{ਜਾਵੇਗੀ।}}
  \item \emph{Read it:} \textbf{\pa{ਕੀ} \pa{ਅਮਨ} \pa{ਕੱਲ੍ਹ} \pa{ਆਵੇਗਾ}?} Give its meaning.
  \item \emph{Write it:} copy \textbf{\pa{ਨਹੀਂ} \pa{ਕਰਾਂਗਾ}} once
\end{itemize}

\begin{tcolorbox}[breakable,colback=teal!4,colframe=teal!35!black,title={Before you move on}]
Ask \emph{will you work tomorrow?} and answer it truthfully.
\end{tcolorbox}
\section[\texorpdfstring{\pa{ਤੁਸੀਂ} \pa{ਕੱਲ੍ਹ} \pa{ਕੀ} \pa{ਕਰੋਗੇ}? (tusīṁ kallh kī karoge?)}{ਤੁਸੀਂ ਕੱਲ੍ਹ ਕੀ ਕਰੋਗੇ? (tusīṁ kallh kī karoge?)}]{Tomorrow's work: a four-skill payoff}

\label{lesson:PA-C222-payoff}



\begin{quote}
Ask \emph{what will you do tomorrow?}
\end{quote}

\subsection*{Your turn: four skills}
\begin{itemize}
\raggedright
  \item \emph{Hear:} \emph{I will work tomorrow} Say it with the ending that fits you.
  \item \emph{Hear:} \emph{will they work tomorrow?} Say \textbf{\pa{ਕੀ} \pa{ਉਹ} \pa{ਕੱਲ੍ਹ} \pa{ਕੰਮ} \pa{ਕਰਨਗੇ}?}
  \item \emph{Read it:} \textbf{\pa{ਨਹੀਂ} \pa{ਜੀ}, \pa{ਅਮਨ} \pa{ਕੱਲ੍ਹ} \pa{ਨਹੀਂ} \pa{ਆਵੇਗਾ।}} Give its meaning: no, Aman won't come tomorrow.
  \item \emph{Write it:} cover the line and write \textbf{\pa{ਤੁਸੀਂ} \pa{ਕੱਲ੍ਹ} \pa{ਕੀ} \pa{ਕਰੋਗੇ}?}
\end{itemize}

\begin{tcolorbox}[breakable,colback=teal!4,colframe=teal!35!black,title={Before you move on}]
Say one thing you will do tomorrow and one thing you will not.
\end{tcolorbox}
